\section*{Insight Từ Hình (Bản Nháp Báo Cáo Nhanh Cho Giáo Sư)}

Phần này chỉ dùng số liệu đã xuất từ pipeline phân tích (không suy diễn ngoài output). Mỗi hình được trình bày theo khung: \textbf{(i) Output quan sát được} và \textbf{(ii) Insight rút ra cho bản thảo}.

\subsection*{Hình 1: Hiệu Ứng Gộp (Kết Quả Cốt Lõi)}
\textbf{Output quan sát được (từ file report).}
\begin{itemize}[leftmargin=1.2em]
  \item Human Augmentation là hiệu ứng dương mạnh nhất: \(g=0.563\), 95\% CI \([0.443, 0.682]\), \(p<0.001\).
  \item Strong Synergy là hiệu ứng âm mạnh nhất: \(g=-0.448\), 95\% CI \([-0.664, -0.233]\), \(p<0.001\).
  \item AI Augmentation có dị hợp cao nhất: \(I^2=99.4\%\).
\end{itemize}
\begin{figure}[H]
  \centering
  \includegraphics[width=\textwidth]{EXP_Figure1_Pooled_Effects.png}
  \caption{Hình 1. Hiệu ứng gộp cho ba đối sánh chính (nguồn: Experiment tự sinh).}
\end{figure}
\textbf{Insight cho bản thảo.} Kết quả gộp cho thấy lợi ích cộng tác bất đối xứng: chiều AI hỗ trợ người là tín hiệu ổn định nhất, trong khi các chiều còn lại biến thiên lớn theo bối cảnh.
\newline
\textit{Nguồn đối chiếu:} \texttt{Experiment/outputs/reports/RESEARCH\_KEY\_FINDINGS\_VI.md}.

\subsection*{Hình 2: Sai Lệch Công Bố (Kiểm Tra Độ Tin Cậy)}
\textbf{Output quan sát được (từ file report).}
\begin{itemize}[leftmargin=1.2em]
  \item Egger có ý nghĩa ở cả 3 đối sánh: \(p=5.706\times10^{-15}\) (es\_s), \(p=0.0002914\) (es\_h), \(p=4.577\times10^{-8}\) (es\_a).
  \item Begg có ý nghĩa ở cả 3 đối sánh: \(p=4.335\times10^{-6}\) (es\_s), \(p=4.362\times10^{-6}\) (es\_h), \(p=1.797\times10^{-5}\) (es\_a).
  \item Trim-and-fill ước lượng số nghiên cứu thiếu: 29 (es\_s), 11 (es\_h), 14 (es\_a).
\end{itemize}
\begin{figure}[H]
  \centering
  \includegraphics[width=\textwidth]{EXP_Figure2_Bias_Diagnostics.png}
  \caption{Hình 2. Chẩn đoán sai lệch công bố: cường độ tín hiệu và số nghiên cứu thiếu.}
\end{figure}
\textbf{Insight cho bản thảo.} Bằng chứng sai lệch công bố nhất quán ở cả ba đối sánh, do đó cần diễn giải thận trọng các hiệu ứng thuận lợi trong tài liệu đã công bố.
\newline
\textit{Nguồn đối chiếu:} \texttt{Experiment/outputs/reports/RESEARCH\_KEY\_FINDINGS\_VI.md}, \texttt{Experiment/outputs/reports/EXPERIMENT\_FIGURE\_INSIGHTS\_VI.md}.

\subsection*{Hình 3: Forest Theo Biến Điều Tiết (Vì Sao Hiệu Ứng Khác Nhau)}
\textbf{Output quan sát được (từ file report).}
\begin{itemize}[leftmargin=1.2em]
  \item Strong Synergy: cao nhất ở Task\_Type=Create (\(g=0.372\), \(k=26\)); thấp nhất ở AI\_Type\_Cleaned=Simulated-AI (\(g=-1.290\), \(k=12\)).
  \item Human Augmentation: cao nhất ở AI\_Type\_Cleaned=Wizard of Oz (\(g=0.734\), \(k=36\)); thấp nhất ở Industry=Communication (\(g=0.276\), \(k=86\)).
  \item AI Augmentation: cao nhất ở Task\_Type=Create (\(g=1.114\), \(k=26\)); thấp nhất ở AI\_Type\_Cleaned=Simulated-AI (\(g=-1.272\), \(k=12\)).
  \item Các moderator có ý nghĩa lặp lại: Task Type, AI Type, Participant Expertise, Industry (và AI\_Expl\_Incl với es\_s/es\_a).
\end{itemize}
\begin{figure}[H]
  \centering
  \includegraphics[width=\textwidth]{EXP_Figure3_Subgroup_Range.png}
  \caption{Hình 3. Biên độ dao động theo điều tiết (cao nhất-thấp nhất).}
\end{figure}
\textbf{Insight cho bản thảo.} Hiệu ứng thay đổi mạnh giữa các nhóm điều tiết và có đảo dấu ở một số cấu hình, cho thấy kết luận theo một mức trung bình toàn cục là không đủ.
\newline
\textit{Nguồn đối chiếu:} \texttt{Experiment/outputs/reports/RESEARCH\_KEY\_FINDINGS\_VI.md}, \texttt{Experiment/outputs/reports/DEEP\_INSIGHTS\_VI.md}.

\subsection*{Hình 4: Hiệu Ứng Theo Ngành (Tính Khái Quát Ngoại Tại)}
\textbf{Output quan sát được (từ file report).}
\begin{itemize}[leftmargin=1.2em]
  \item Strong Synergy: tốt nhất Healthcare (\(g=-0.306\)), thấp nhất Business (\(g=-0.926\)); \(Q\)-between \(p<0.001\).
  \item Human Augmentation: cao nhất Business (\(g=0.706\)), thấp nhất Communication (\(g=0.276\)); \(Q\)-between \(p<0.001\).
  \item AI Augmentation: cao nhất Communication (\(g=1.011\)), thấp nhất Business (\(g=-0.732\)); \(Q\)-between \(p<0.001\).
\end{itemize}
\begin{figure}[H]
  \centering
  \includegraphics[width=\textwidth]{EXP_Figure4_Industry_Profile.png}
  \caption{Hình 4. Cấu hình hiệu ứng theo ngành cho ba đối sánh chính.}
\end{figure}
\textbf{Insight cho bản thảo.} Industry là điều kiện biên quan trọng: cùng một chiến lược cộng tác có thể đổi độ lớn, thậm chí đổi dấu hiệu ứng giữa các miền.
\newline
\textit{Nguồn đối chiếu:} \texttt{Experiment/outputs/reports/RESEARCH\_KEY\_FINDINGS\_VI.md}.

\subsection*{Tóm Tắt 1 Đoạn Cho Slide/Meeting}
Từ các output tự động, có thể chốt ba điểm: (1) hiệu ứng gộp bất đối xứng, nổi bật nhất là Human Augmentation dương; (2) sai lệch công bố xuất hiện nhất quán ở cả ba đối sánh; (3) điều tiết theo thiết kế nhiệm vụ/loại AI/ngành tạo khác biệt lớn, gồm cả đảo dấu hiệu ứng. Bản thảo nên vì vậy nhấn mạnh tính phụ thuộc bối cảnh của cộng tác người--AI thay vì diễn giải theo một hiệu ứng chung.
